%!TEX root = ../main.tex
\section{Introduction}
\label{sec:introduction}

The model that writes the best scientific figure descriptions is also the one that most confidently fabricates information it cannot see.
Vision-language models are rapidly entering scientific workflows, from literature review assistants that summarize chart findings to accessibility tools that narrate figures for visually impaired researchers.
When a model misreads a trend line or invents a data point, the downstream consequences propagate silently: a fabricated value becomes a cited fact, a hallucinated comparison shapes a research conclusion.
The stakes demand not only that models describe figures accurately, but that they behave reliably when faced with ambiguity, missing information, or misleading context.
Yet current evaluation practices assess only the first of these requirements, leaving the second entirely unmeasured.

Existing benchmarks for scientific figure understanding evaluate models exclusively through question answering on clean images.
FigureQA~\citep{kahou2018figureqa} and PlotQA~\citep{methani2020plotqa} test binary and numerical reasoning over synthetic or semi-synthetic charts.
ChartQA~\citep{masry2022chartqa} advances to human-authored questions on real-world figures, and ChartBench~\citep{xu2024chartbench} broadens chart type coverage.
These benchmarks measure whether a model produces a correct answer, but none ask whether a model admits when it cannot see, resists false premises embedded in questions, or distinguishes its own visual evidence from authoritative-sounding but incorrect textual context.
A model can score highly on ChartQA while confidently fabricating information about chart elements that do not exist, because no existing benchmark tests for this failure mode.

We present \textsc{SciFig-Eval}, a diagnostic evaluation framework that measures both description quality and behavioral reliability on scientific figures.
The framework is organized around two evaluation axes (open-ended description and targeted reasoning) crossed with two conditions (standard and adversarial), yielding a $2\!\times\!2$ evaluation matrix that structures every experiment.
Adversarial probes are grounded in cognitive science: \emph{presupposition embedding}~\citep{loftus1975leading} tests whether definite articles trick models into accepting non-existent chart elements, \emph{anchoring bias}~\citep{tversky1974judgment} plants plausible but incorrect values, and \emph{sycophantic agreement}~\citep{sharma2024sycophancy} measures whether models echo false captions over their own visual perception.
These probes feed into the \emph{A-R-I behavioral framework}, which decomposes model behavior into three independently measurable dimensions: Admittance (does the model acknowledge visual limitations?), Resistance (does it push back on false premises?), and Inductance (can it infer missing information from context?).
We evaluate 8 models across 250 figures and over 11,000 evaluation instances (Figure~\ref{fig:overview}).

Quality rankings diverge from reliability rankings at precisely the points that matter for deployment.
GPT-5.2 leads on description quality with an MQM score of 91.6, but admits visual limitations only 6\% of the time when asked about elements it cannot see, fabricating answers for the remaining 94\%.
Gemini 3.1 Pro, ranking second on quality at 90.2, admits uncertainty 90\% of the time and achieves the highest resistance to hallucination probes (0.91).
The divergence extends to caption dependency: Phi-4 follows poisoned captions over its own visual evidence 95\% of the time ($R = 0.05$), while presupposition embedding, the technique borrowed from eyewitness testimony research, catches models that successfully resist explicit false values but fail against implicit assumptions about non-existent elements.
A benchmark reporting only quality scores would rank GPT-5.2 and Gemini as near-equivalent, missing the fact that one fabricates 98\% of the time while the other acknowledges uncertainty.

Our contributions are as follows:
\begin{enumerate}
    \item A multi-dimensional evaluation framework for scientific figure understanding that pairs standard quality assessment with psychology-informed adversarial probes across two evaluation axes and two conditions (\S\ref{sec:framework}).
    \item The A-R-I behavioral framework, which decomposes model behavior into Admittance, Resistance, and Inductance, providing independently measurable dimensions of reliability under uncertainty (\S\ref{sec:ari}).
    \item An empirical demonstration that description quality and behavioral reliability are distinct dimensions of VLM competence, with model rankings diverging substantially across these dimensions (\S\ref{sec:results}).
    \item A full benchmark release including the annotated dataset, adversarial probes, evaluation code, and an interactive review dashboard.
\end{enumerate}
